\documentclass[11pt]{article}

% --- Idioma y fuentes ------------------------------------------------
\usepackage[spanish,es-noshorthands,es-tabla]{babel}
\usepackage[T1]{fontenc}
\usepackage{lmodern}

% --- Geometría y encabezado ------------------------------------------
\usepackage[a4paper,margin=2.5cm]{geometry}
\usepackage{fancyhdr}
\setlength{\headheight}{38.22325pt}
\addtolength{\topmargin}{-6.22325pt}

% --- Matemáticas ------------------------------------------------------
\usepackage{amsmath,amssymb}

% --- Colores, tablas, listas, código ----------------------------------
\usepackage[dvipsnames,table]{xcolor}
\usepackage{booktabs}
\usepackage{array}
\usepackage{longtable}
\usepackage{enumitem}
\usepackage{listings}
\renewcommand{\arraystretch}{1.5}

\lstset{
  language=R,
  basicstyle=\ttfamily\footnotesize,
  commentstyle=\color{gray}\itshape,
  keywordstyle=\color{blue!70!black},
  stringstyle=\color{teal!60!black},
  breaklines=true,
  columns=fullflexible,
  showstringspaces=false
}


\pagestyle{fancy}
\fancyhf{}
\fancyhead[L]{\textbf{Econometría II}\\ Trabajo práctico --- Etapa 1}
\fancyhead[C]{Universidad de Santiago de Chile}
\fancyfoot[C]{\thepage}

\begin{document}
\sloppy

% =====================================================================
%  PORTADA
% =====================================================================
\begin{titlepage}
\centering
\vspace*{3cm}
{\large Universidad de Santiago de Chile}\\[6pt]
{\large Econometría II}\\[2cm]
{\LARGE \textbf{Etapa 1: Construcción de la base de datos}}\\[10pt]
{\large Trabajo Práctico --- Admisión Universidad 2026}\\[3cm]

\textbf{Integrantes:}\quad
David Cuevas $\bullet$ Javier Li-Yau $\bullet$ Benjamín Lizana $\bullet$ Ana Nacif

\vspace{0.6cm}

\textbf{Profesor:}\quad Brian Pfeng Ibacache
\qquad
\textbf{Fecha:}\quad 8 de octubre de 2026
\end{titlepage}

% =====================================================================
\subsection*{1 \; Introducción}
% =====================================================================

Con cuatro archivos del DEMRE del proceso de admisión 2026 armamos una
base de 60.544 estudiantes que egresaron de enseñanza media en 2025 y
que se matricularon por vía regular, con tipo de matrícula postulación
regular, a lo más en su tercera preferencia. Ese filtro deja a todos
dentro de sus tres primeras opciones, así que por sí solo no permite
distinguirlos. Lo que sí cambia de uno a otro es en cuál de las tres
quedaron: 40.173 (66,4\,\%) lograron la primera y 20.371 (33,6\,\%) se
matricularon en la segunda o la tercera. Nuestro trabajo apunta a
explicar de qué depende quedar en la primera preferencia, y esa división
es la variable de respuesta, \texttt{PREF\_PRIMERA\_DUMMY}.

La base tiene una fila por estudiante y 29 variables. De cada uno trae
su promedio de notas, sus puntajes NEM y ranking, sus puntajes de
Competencia Lectora y Matemática 1, su sexo, su tramo de ingreso
familiar, si fue preseleccionado PACE o BEA y la preferencia en que se
matriculó. De su colegio de egreso trae la dependencia, la rama
educacional, la región y la comuna, y tres variables que calculamos
nosotros: inscritos en la PAES, matriculados y promedio del tramo de
ingreso de sus inscritos. Con esta información se estimará el modelo
logit de la Etapa 3.

\medskip\hrule\medskip

% =====================================================================
\subsection*{2 \; Fuentes de datos}
% =====================================================================

Todos los datos vienen del portal de bases de datos del DEMRE, proceso de
admisión 2026. Usamos cuatro archivos, que se unen por el identificador
del estudiante, \texttt{ID\_aux}.

\begin{center}
\begin{tabular}{>{\raggedright\arraybackslash}p{3.4cm}
                >{\raggedright\arraybackslash}p{6.4cm}
                >{\raggedleft\arraybackslash}p{2.6cm}}
\toprule
\textbf{Archivo} & \textbf{Contenido} & \textbf{Filas} \\
\midrule
B (\texttt{df1}) & Inscritos: colegio, domicilio, sexo, ingreso familiar & 317.700 \\
C (\texttt{df2}) & Notas, NEM, ranking y puntajes PAES & 317.700 \\
D (\texttt{df3}) & Postulaciones (una fila por carrera postulada) & 1.686.689 \\
Matrícula (\texttt{df4}) & Vía, tipo de matrícula, preferencia y puntaje & 128.345 \\
\bottomrule
\end{tabular}
\end{center}


\medskip\hrule\medskip

% =====================================================================
\subsection*{3 \; Construcción de la base}
% =====================================================================

\subsubsection*{3.1 \; Variables del colegio}

Antes de filtrar, calculamos las variables del colegio agrupando por
establecimiento (\texttt{RBD}). Se calculan antes porque necesitan a todos
los estudiantes del colegio, no solo a los que quedan en la base final.

\begin{itemize}[leftmargin=1.5em, itemsep=0.3em]
  \item Del archivo B: número de inscritos del colegio y promedio del
        ingreso per cápita familiar de sus inscritos.
  \item Del archivo de matrícula: número de matriculados por vía regular y
        tipo de matrícula 1 que egresaron del colegio. Los colegios sin
        matriculados quedan con 0.
\end{itemize}

\subsubsection*{3.2 \; Selección de columnas y unión}

De cada archivo dejamos solo las columnas que usamos (detalle en la
sección 7). El archivo D trae una fila por carrera postulada; nos quedamos con las
postulaciones de tipo regular y lo dejamos en una fila por estudiante
(188.538 estudiantes). El archivo de
matrícula ya trae una fila por estudiante.

Luego unimos los cuatro archivos por \texttt{ID\_aux} con
\texttt{full\_join}, lo que da 317.700 filas, y a esa base le aplicamos
los filtros.

\subsubsection*{3.3 \; Filtros aplicados}

\begin{enumerate}[leftmargin=1.5em, itemsep=0.3em, label=\textbf{\alph*)}]
  \item \textbf{Promoción del año:} estudiantes que egresaron de enseñanza
        media en 2025 (\texttt{ANYO\_EGRESO} igual a 2025).
  \item \textbf{Vía y tipo de matrícula:} matriculados por vía regular
        (\texttt{VIA} igual a 1) con tipo de matrícula «postulación
        regular» (\texttt{TIPO\_MATRICULA} igual a 1).
  \item \textbf{Inscripción en el proceso:} inscribieron al menos una
        prueba del proceso 2026 (\texttt{RINDIO\_PROCESO\_ACTUAL} igual
        a 1).
  \item \textbf{Preferencia:} matriculados en su primera, segunda o
        tercera preferencia (\texttt{PREFERENCIA} entre 1 y 3).
\end{enumerate}

Además, sacamos a quienes no tienen puntaje NEM o lo tienen igual a 0,
porque sin notas de enseñanza media no se pueden comparar con el resto.

\subsubsection*{3.4 \; Observaciones en cada paso}

\begin{center}
\begin{tabular}{>{\raggedright\arraybackslash}p{9cm}
                >{\raggedleft\arraybackslash}p{3cm}}
\toprule
\textbf{Paso} & \textbf{Observaciones} \\
\midrule
Unión de los cuatro archivos & 317.700 \\
Egresados en 2025 & 218.989 \\
Vía regular y tipo de matrícula 1 & 69.004 \\
Inscribieron alguna prueba del proceso 2026 & 69.002 \\
Matriculados en preferencia 1 a 3 & 60.558 \\
\rowcolor{yellow!30}
\textbf{Con puntaje NEM distinto de 0 (base final)} & \textbf{60.544} \\
\bottomrule
\end{tabular}
\end{center}

\medskip\hrule\medskip

% =====================================================================
\subsection*{4 \; Variable de respuesta}
% =====================================================================

Por los filtros, todos los estudiantes de la base se matricularon en una
de sus tres primeras preferencias. Por eso la variable de respuesta separa
a quienes quedaron en su primera preferencia de quienes quedaron en la
segunda o la tercera:

\[
\texttt{PREF\_PRIMERA\_DUMMY}_i =
\begin{cases}
1 & \text{si se matriculó en su primera preferencia} \\[4pt]
0 & \text{si se matriculó en su segunda o tercera preferencia}
\end{cases}
\]

\begin{center}
\begin{tabular}{lrr}
\toprule
\textbf{Valor} & \textbf{Estudiantes} & \textbf{Porcentaje} \\
\midrule
1 (primera preferencia) & 40.173 & 66,4\% \\
0 (segunda o tercera)   & 20.371 & 33,6\% \\
\midrule
\textbf{Total} & \textbf{60.544} & 100\% \\
\bottomrule
\end{tabular}
\end{center}

De los 20.371 que no quedaron en su primera preferencia, 13.322 se
matricularon en la segunda y 7.049 en la tercera.

\medskip\hrule\medskip

% =====================================================================
\subsection*{5 \; Variables creadas}
% =====================================================================

\subsubsection*{5.1 \; Variables del establecimiento}

Se pegan a cada estudiante por el \texttt{RBD} de su colegio de egreso.
Así cada estudiante tiene información de su colegio en el mismo proceso
2026.

\begin{itemize}[leftmargin=1.5em, itemsep=0.3em]
  \item \texttt{N\_ALUMNOS}: número de inscritos en el proceso 2026 que
        egresaron del colegio.
  \item \texttt{INGRESO\_COLEGIO}: promedio de
        \texttt{INGRESO\_PERCAPITA\_GRUPO\_FA} de los inscritos del
        colegio. El promedio incluye el código 99 («prefiero no
        responder»), así que no queda en la escala de 1 a 10 de los
        tramos: su promedio es 31,8 (sección 8).
  \item \texttt{N\_MATRICULADOS}: número de egresados del colegio que se
        matricularon por vía regular y tipo de matrícula 1.
\end{itemize}

Los 244 estudiantes sin colegio de egreso registrado (sección 6) quedan
agrupados en el cálculo como si fueran un solo colegio, con 3.169
inscritos y 902 matriculados. Para ellos estas tres variables no
describen un establecimiento real. Sin ese grupo, el colegio más grande
tiene 947 inscritos y 447 matriculados.

No incluimos si el colegio es rural ni el número de alumnos por curso,
porque las bases del DEMRE no traen esa información: para construirlas
se necesitan los registros de establecimientos del MINEDUC.

\subsubsection*{5.2 \; Variables dummy}

\begin{itemize}[leftmargin=1.5em, itemsep=0.3em]
  \item \texttt{PREF\_PRIMERA\_DUMMY}: variable de respuesta (sección 4).
  \item \texttt{TP\_DUMMY}: 1 si el colegio es técnico profesional
        (rama T1 a T5), 0 si no. Hay 9.110 estudiantes de colegios
        técnico profesionales y 51.434 del resto. Entre esos 51.434 están los
        244 sin colegio registrado, que quedan con 0 aunque no se conoce
        su rama.
\end{itemize}

\medskip\hrule\medskip

% =====================================================================
\subsection*{6 \; Datos sin información}
% =====================================================================

\begin{center}
\begin{tabular}{>{\raggedright\arraybackslash}p{5.4cm}
                >{\raggedleft\arraybackslash}p{1.8cm}
                >{\raggedright\arraybackslash}p{5.6cm}}
\toprule
\textbf{Variable} & \textbf{Casos} & \textbf{Qué significa} \\
\midrule
\texttt{RBD}, \texttt{RAMA\_EDUCACIONAL}, \texttt{GRUPO\_DEPENDENCIA},
\texttt{CODIGO\_COMUNA} & 244 & Estudiante sin colegio de egreso registrado. \\
\texttt{CODIGO\_REGION} & 10 & Región del colegio sin dato. \\
\texttt{PACE} & 55.382 & No fue preseleccionado PACE. \\
\texttt{BEA} & 53.805 & No fue preseleccionado a la Beca de Excelencia Académica. \\
\texttt{INGRESO\_PERCAPITA\_GRUPO\_FA} igual a 99 & 16.235 & «Prefiero no responder». \\
\bottomrule
\end{tabular}
\end{center}

En el resto de las variables no hay datos faltantes. Para los 244
estudiantes sin colegio, sin embargo, \texttt{N\_ALUMNOS},
\texttt{INGRESO\_COLEGIO}, \texttt{N\_MATRICULADOS} y \texttt{TP\_DUMMY}
tienen valores que no corresponden a un colegio (sección 5).

\medskip\hrule\medskip

% =====================================================================
\subsection*{7 \; Definición de las variables}
% =====================================================================

Fuente: B, C, D y M (matrícula) son los archivos del DEMRE; «creada»
indica que la construimos nosotros.

\begin{longtable}{>{\raggedright\arraybackslash}p{5.4cm}
                  >{\raggedright\arraybackslash}p{1.4cm}
                  >{\raggedright\arraybackslash}p{6.6cm}}
\toprule
\textbf{Variable} & \textbf{Fuente} & \textbf{Definición} \\
\midrule
\endfirsthead
\toprule
\textbf{Variable} & \textbf{Fuente} & \textbf{Definición} \\
\midrule
\endhead
\bottomrule
\endfoot
\multicolumn{3}{l}{\textit{Identificación y respuesta}} \\
\texttt{ID\_aux} & B, C, D, M & Identificador único del estudiante. \\
\texttt{PREF\_PRIMERA\_DUMMY} & Creada & Variable de respuesta: 1 si se matriculó en su primera preferencia, 0 si en la segunda o tercera. \\
\midrule
\multicolumn{3}{l}{\textit{Estudiante}} \\
\texttt{FECHA\_NACIMIENTO} & B & Mes y año de nacimiento (MMAAAA). Es numérica, así que los meses 1 a 9 quedan sin el cero inicial: 42006 es abril de 2006. \\
\texttt{SEXO} & B & 1 hombre, 2 mujer. \\
\texttt{CODIGO\_COMUNA\_D} & B & Comuna de domicilio del estudiante. \\
\texttt{INGRESO\_PERCAPITA\_GRUPO\_FA} & B & Tramo de ingreso per cápita del grupo familiar, de 1 a 10; 99 es «prefiero no responder». \\
\texttt{PACE} & B & Preseleccionado PACE por el MINEDUC. \\
\texttt{BEA} & B & Preseleccionado a la Beca de Excelencia Académica. \\
\texttt{ANYO\_EGRESO} & B & Año de egreso de enseñanza media (2025 en toda la base). \\
\texttt{RINDIO\_PROCESO\_ACTUAL} & B & 1 si inscribió al menos una prueba del proceso 2026 (1 en toda la base). \\
\midrule
\multicolumn{3}{l}{\textit{Notas y puntajes}} \\
\texttt{PROMEDIO\_NOTAS} & C & Promedio de notas de enseñanza media. \\
\texttt{PTJE\_NEM} & C & Puntaje NEM (notas de enseñanza media). \\
\texttt{PTJE\_RANKING} & C & Puntaje ranking de notas. \\
\texttt{CLEC\_REG\_ACTUAL} & C & Puntaje PAES de Competencia Lectora, rendición regular 2026. \\
\texttt{MATE1\_REG\_ACTUAL} & C & Puntaje PAES de Competencia Matemática 1, rendición regular 2026. \\
\midrule
\multicolumn{3}{l}{\textit{Postulación y matrícula}} \\
\texttt{TIPO\_PREF} & D & Tipo de postulación (REGULAR en toda la base). \\
\texttt{PREFERENCIA} & M & Preferencia en que se matriculó (1 a 3 en la base). Con ella se construye la variable de respuesta. \\
\texttt{VIA} & M & Vía de matrícula (1, regular, en toda la base). \\
\texttt{TIPO\_MATRICULA} & M & Tipo de matrícula (1, postulación regular, en toda la base). \\
\texttt{PTJE\_POND} & M & Puntaje ponderado con que se matriculó. \\
\midrule
\multicolumn{3}{l}{\textit{Colegio de egreso}} \\
\texttt{RBD} & B & Identificador del establecimiento. \\
\texttt{RAMA\_EDUCACIONAL} & B & H1 y H2 humanista científico (diurno y nocturno); T1 a T5 técnico profesional. \\
\texttt{GRUPO\_DEPENDENCIA} & B & 1 particular pagado, 2 particular subvencionado, 3 municipal, 4 servicio local de educación. \\
\texttt{CODIGO\_REGION} & B & Región del establecimiento. \\
\texttt{CODIGO\_COMUNA} & B & Comuna del establecimiento. \\
\texttt{TP\_DUMMY} & Creada & 1 si el colegio es técnico profesional (T1 a T5), 0 si no. \\
\texttt{N\_ALUMNOS} & Creada & Inscritos del proceso que egresaron del colegio. \\
\texttt{INGRESO\_COLEGIO} & Creada & Promedio de \texttt{INGRESO\_PERCAPITA\_GRUPO\_FA} de los inscritos del colegio, incluido el código 99. \\
\texttt{N\_MATRICULADOS} & Creada & Egresados del colegio matriculados por vía regular y tipo de matrícula 1. \\
\end{longtable}

\texttt{PREFERENCIA} y \texttt{PTJE\_POND} quedan en la base como
información de la matrícula, pero no se usarán como variables
explicativas: la variable de respuesta se construye con
\texttt{PREFERENCIA}, y \texttt{PTJE\_POND} es el puntaje ponderado de
la carrera en que el estudiante ya quedó matriculado, así que depende
del resultado.

\medskip\hrule\medskip

% =====================================================================
\subsection*{8 \; Base final}
% =====================================================================

La base final (\texttt{df\_final}) tiene \textbf{60.544 filas y 29
columnas}, con una fila por estudiante: cada \texttt{ID\_aux} aparece una
sola vez. Los estudiantes vienen de 3.108 colegios distintos.

\begin{center}
\begin{tabular}{lrrrr}
\toprule
\textbf{Variable} & \textbf{Mínimo} & \textbf{Mediana} & \textbf{Promedio} & \textbf{Máximo} \\
\midrule
\texttt{PROMEDIO\_NOTAS}   & 4,50  & 6,33  & 6,27  & 7,00 \\
\texttt{PTJE\_NEM}         & 259   & 808   & 792   & 1.000 \\
\texttt{PTJE\_RANKING}     & 259   & 822   & 806   & 1.000 \\
\texttt{CLEC\_REG\_ACTUAL} & 316   & 663   & 667   & 1.000 \\
\texttt{MATE1\_REG\_ACTUAL}& 320   & 682   & 698   & 1.000 \\
\texttt{PTJE\_POND}        & 354,6 & 716,1 & 718,6 & 998,4 \\
\texttt{N\_ALUMNOS}        & 8     & 127   & 172   & 3.169 \\
\texttt{INGRESO\_COLEGIO}  & 1,0   & 29,7  & 31,8  & 94,1 \\
\texttt{N\_MATRICULADOS}   & 1     & 55    & 74    & 902 \\
\bottomrule
\end{tabular}
\end{center}

En la base hay 31.981 mujeres (52,8\%) y 28.563 hombres (47,2\%).

\medskip\hrule\medskip

% =====================================================================
\subsection*{9 \; Conclusión}
% =====================================================================

La base final tiene 60.544 estudiantes, uno por fila: el 66,4\,\% quedó
en su primera preferencia y el 33,6\,\% en la segunda o la tercera. Como
todos cumplen los mismos filtros, esa diferencia se estudiará con sus
notas, sus puntajes, su ingreso familiar y las variables de su colegio.
\texttt{PREFERENCIA} y \texttt{PTJE\_POND} quedan fuera del modelo porque
dependen del resultado. Antes de estimar hay que recalcular
\texttt{INGRESO\_COLEGIO} sin el código 99 y dejar sin valor las
variables del colegio de los 244 estudiantes sin \texttt{RBD}.

% =====================================================================
\end{document}
